% ============================================================
%  CAPÍTULO 4 — ARQUITECTURA NEURONAL PROPUESTA
%  Memoria TFG: Transcripción automática de música de piano
%  Autor: María Piedad Castex Sierra
%  Formato: XeLaTeX — mismo estilo que Capítulos 3 y 4
% ============================================================

\chapter{Arquitectura Neuronal Propuesta}
\label{cap:arquitectura}

El presente capítulo describe en detalle la arquitectura neuronal diseñada e implementada
para el sistema de transcripción automática de piano. Partiendo de los fundamentos
teóricos establecidos en el Capítulo~\ref{cap:estado_del_arte} y de las representaciones
definidas en el Capítulo~\ref{cap:datos}, se describe la configuración global del modelo,
la estructura del codificador basado en atención dispersa, el mecanismo de predicción de
máscara que constituye el núcleo de la contribución técnica de este trabajo, y el
decodificador autorregresivo que genera la secuencia de \textit{tokens} MIDI. Finalmente,
se detalla el esquema de codificación posicional adoptado y se justifican las decisiones de
diseño más relevantes.

% -----------------------------------------------------------
\section{Visión general del modelo}
\label{sec:vision_general}

Siguiendo el paradigma de traducción de secuencia a secuencia
descrito en la Sección~\ref{sec:seq2seq}, el sistema establece
un mapeo directo entre la representación acústica de la pieza y
la secuencia simbólica de eventos MIDI. El sistema desarrollado
en este proyecto, denominado \texttt{Magdalena v1.0}, adopta una
arquitectura Encoder-Decoder asimétrica: el codificador y el
decodificador difieren tanto en su estructura interna como en el
mecanismo de atención que emplean, lo que distingue esta propuesta
de las implementaciones estándar del Transformer original
\cite{vaswani2017attention}.

El flujo de procesamiento del modelo puede describirse formalmente como sigue. Dado un
espectrograma Log-Mel de entrada $\mathbf{X} \in \mathbb{R}^{B \times N_{mels} \times T}$,
donde $B$ es el tamaño del lote (\textit{batch size}), $N_{mels} = 229$ es el número de
bandas frecuenciales y $T$ es la longitud temporal de la secuencia, el codificador produce
una representación latente contextualizada $\mathbf{M} \in \mathbb{R}^{B \times T \times d}$,
denominada memoria (\textit{memory}), con dimensión de embedding $d = 256$. El
decodificador consume esta memoria y genera, de forma autorregresiva, la secuencia de logits $\mathbf{L} \in \mathbb{R}^{B \times S \times |\mathcal{V}|}$,
donde $S$ es la longitud de la secuencia objetivo y $|\mathcal{V}| = 311$ es el tamaño
del vocabulario definido en la Sección~\ref{sec:dominio_simbolico}. La configuración
Encoder-Decoder y el mecanismo de autoatención subyacente fueron descritos en detalle
en la Sección~\ref{sec:transformer}.

La Tabla~\ref{tab:hiperparametros_globales} recoge los hiperparámetros
del modelo tal y como están configurados en los experimentos realizados.

\begin{table}[H]
    \centering
    \renewcommand{\arraystretch}{1.3}
    \begin{tabular}{l l p{7cm}}
        \hline
        \textbf{Hiperparámetro} & \textbf{Valor} & \textbf{Descripción} \\
        \hline
        $d$ (\texttt{embed\_dim})
            & 256\textsuperscript{†}
            & Dimensión del espacio de \textit{embedding} en todas
              las capas. \\
        $H$ (\texttt{nhead})
            & 4\textsuperscript{†}
            & Número de cabezas de atención (\textit{multi-head}). \\
        $L_{enc}$ (\texttt{num\_encoder\_layers})
            & 4\textsuperscript{†}
            & Profundidad del codificador (bloques Sparsifiner). \\
        $L_{dec}$ (\texttt{num\_decoder\_layers})
            & 4\textsuperscript{†}
            & Profundidad del decodificador (bloques Transformer
              estándar). \\
        $d_{ff}$ (\texttt{dim\_feedforward})
            & 2048
            & Dimensión de la red \textit{feed-forward} interna
              ($8 \times d$). \\
        $p$ (\texttt{dropout})
            & 0.1
            & Tasa de \textit{dropout} aplicada en
              \textit{embeddings} y atención. \\
        $T_{train}$ (\texttt{max\_audio\_frames})
            & 2\,048 -- 4\,096
            & Longitud máxima de secuencia de audio en
              entrenamiento (varía por experimento, véase
              Capítulo~\ref{cap:implementacion}); capacidad máxima
              del modelo: 70\,000 tramas. \\
        $r_{attn}$ (\texttt{attn\_keep\_rate})
            & 0.25
            & Fracción de conexiones de atención retenidas por
              Sparsifiner. \\
        $r_{N}$ (\texttt{reduce\_n\_factor})
            & 64
            & Factor de compresión temporal en el predictor de
              máscara. \\
        \textit{DropPath} (\texttt{drop\_path\_rate})
            & 0.0
            & Regularización de profundidad estocástica
              implementada pero configurada como inactiva en los
              experimentos actuales (\texttt{nn.Identity}). \\
        \hline
    \end{tabular}
    \caption{Hiperparámetros del modelo \texttt{Magdalena v1.0}
    en la configuración experimental. \textsuperscript{†}El código
    soporta una configuración de mayor capacidad
    ($d = 512$, $H = 8$, $L_{enc} = L_{dec} = 6$), equivalente
    al Transformer original \cite{vaswani2017attention}, descartada
    por las limitaciones de memoria de la GPU disponible
    (Google Colab, T4 de 15\,GB de VRAM). La configuración
    reducida mantiene un \textit{batch} efectivo de 32 secuencias
    en todos los experimentos, alcanzado mediante distintas
    combinaciones de tamaño de \textit{batch} y acumulación de
    gradientes según las restricciones de memoria de cada
    configuración (véase Capítulo~\ref{cap:implementacion}).}
    \label{tab:hiperparametros_globales}
\end{table}

% -----------------------------------------------------------
\section{Codificador con atención dispersa (Sparsifiner)}
\label{sec:encoder_sparsifiner}
% -----------------------------------------------------------

El codificador del sistema, implementado en la clase \texttt{AudioSparsifinerEncoder},
constituye la principal contribución técnica de esta arquitectura. A diferencia del
codificador estándar del Transformer, que aplica autoatención densa con complejidad
$\mathcal{O}(T^2)$, el codificador propuesto integra el mecanismo Sparsifiner
\cite{wei2023sparsifiner}, que aprende a construir de forma dinámica y dependiente de
la instancia un grafo de conectividad dispersa sobre la secuencia de tramas del
espectrograma.

\subsection{Proyección de entrada y embedding posicional}

El primer componente del codificador es una capa de proyección lineal
$\mathbf{W}_{in} \in \mathbb{R}^{N_{mels} \times d}$ que transforma cada trama
frecuencial del espectrograma —un vector de dimensión $N_{mels} = 229$— en un vector
de embedding de dimensión $d = 256$. Formalmente, dado el espectrograma de entrada
$\mathbf{X} \in \mathbb{R}^{B \times N_{mels} \times T}$, se permutan las dimensiones
frecuencial y temporal para obtener $\mathbf{X}' \in \mathbb{R}^{B \times T \times N_{mels}}$,
y la proyección se aplica sobre la última dimensión:

\begin{equation}
    \mathbf{E} = \mathbf{X}' \, \mathbf{W}_{in} + \mathbf{b}_{in} \in \mathbb{R}^{B \times T \times d}
    \label{eq:proyeccion_entrada}
\end{equation}

A la representación resultante se le suma un embedding posicional aprendido
$\mathbf{P} \in \mathbb{R}^{1 \times T_{max} \times d}$, inicializado con distribución
normal truncada de desviación estándar $\sigma = 0.02$ y recortado dinámicamente a la
longitud real de la secuencia en cada paso de inferencia:

\begin{equation}
    \mathbf{Z}^{(0)} = \mathbf{E} + \mathbf{P}_{[:T, :]} \in \mathbb{R}^{B \times T \times d}
    \label{eq:embedding_posicional_encoder}
\end{equation}

La elección de un embedding posicional aprendido, en lugar del encoding sinusoidal fijo
utilizado en el decodificador, está motivada por la naturaleza bidimensional del
espectrograma: al interpretar las tramas temporales como \textit{tokens} de una
representación matricial, el modelo puede aprender patrones posicionales específicos del
dominio acústico que no están bien capturados por funciones periódicas fijas
\cite{dosovitskiy2021image}.

\subsection{Bloques de atención Sparsifiner}
\label{subsec:bloques_sparsifiner}

El núcleo del codificador está formado por una pila de $L_{enc} = 4$ bloques Sparsifiner
idénticos en estructura, implementados en la clase \texttt{Block}. Cada bloque sigue la
arquitectura Pre-LayerNorm estándar del Transformer \cite{xiong2020layer}, con conexiones
residuales y una red de propagación hacia adelante (\textit{feed-forward network}, FFN)
con función de activación GELU \cite{hendrycks2016gaussian}:

\begin{align}
    \mathbf{Z}'^{(l)} &= \mathbf{Z}^{(l-1)} + \text{DropPath}\!\left(\text{SparseAttn}\!\left(\text{LN}\!\left(\mathbf{Z}^{(l-1)}\right), \boldsymbol{\pi}^{(l-1)}\right)\right) \label{eq:bloque_attn} \\
    \mathbf{Z}^{(l)}  &= \mathbf{Z}'^{(l)} + \text{DropPath}\!\left(\text{FFN}\!\left(\text{LN}\!\left(\mathbf{Z}'^{(l)}\right)\right)\right) \label{eq:bloque_ffn}
\end{align}

donde $\boldsymbol{\pi}^{(l-1)} \in \{0,1\}^{B \times (T-1)}$ es la máscara de
\textit{tokens} activos propagada desde el bloque anterior, $\text{LN}(\cdot)$ denota
normalización de capa (\textit{Layer Normalization}) \cite{ba2016layer}, y $\text{DropPath}(\cdot)$ es la regularización de profundidad
estocástica (\textit{Stochastic Depth}) \cite{huang2016deep},
que descarta bloques residuales completos de forma aleatoria
durante el entrenamiento para prevenir el sobreajuste. En los
experimentos actuales esta regularización se mantiene inactiva
(\texttt{drop\_path\_rate = 0.0}), equivaliendo a la función
identidad; la arquitectura la soporta de forma nativa para
experimentos futuros con modelos más profundos.
La dimensión oculta de la FFN es $d_{ff} = 2048$, lo que supone
una relación $d_{ff} = 8d$ con la configuración experimental
($d = 256$). Esta proporción, mayor que la relación $4d$
convencional del Transformer original \cite{vaswani2017attention},
compensa parcialmente la reducción de $d$ manteniendo la capacidad
expresiva de la red de propagación hacia adelante.

% -----------------------------------------------------------
\section{Mecanismo de predicción de máscara de atención}
\label{sec:mask_predictor}
% -----------------------------------------------------------

El componente diferenciador de la arquitectura es el \texttt{MaskPredictor}, un módulo
auxiliar de bajo coste computacional que, a partir de las representaciones de consulta
(\textit{query}) y clave (\textit{key}) de la capa de atención, predice qué pares de
posiciones son suficientemente relevantes para justificar el cómputo de atención completa.
Este predictor opera en tres fases: compresión de rango, cálculo de atención aproximada
y generación de la máscara binaria final.

\subsection{Compresión de rango}
\label{subsec:compresion_rango}

Para calcular el mapa de relevancia entre posiciones sin incurrir en el coste cuadrático
de la atención completa, el \texttt{MaskPredictor} proyecta las representaciones de
consulta y clave a un espacio de dimensión reducida mediante dos operaciones de
compresión independientes y paralelas.

La primera compresión actúa sobre la dimensión de canales: cada cabeza de atención tiene
dimensión $d_h = d / H = 64$, y se proyecta a una dimensión reducida
$d_c = d_h / r_c = 32$, donde el factor de reducción de canal es $r_c = 2$:

\begin{align}
    \tilde{\mathbf{Q}} &= \mathbf{Q} \, \mathbf{W}_{Qc} \in \mathbb{R}^{B \times H \times T \times d_c} \label{eq:q_reducido} \\
    \tilde{\mathbf{K}} &= \mathbf{K} \, \mathbf{W}_{Kc} \in \mathbb{R}^{B \times H \times T \times d_c} \label{eq:k_reducido}
\end{align}

La segunda compresión actúa sobre la dimensión temporal: la matriz de claves reducida
se proyecta sobre un conjunto de $k = \lfloor T / r_N \rfloor$ vectores base aprendidos,
donde el factor de reducción temporal es $r_N = 64$ (definido en \texttt{transformer.py}).
Esta proyección se realiza mediante una matriz de parámetros aprendidos
$\mathbf{P}_N \in \mathbb{R}^{T_{max} \times k}$, recortada dinámicamente a la longitud
real de la secuencia $T$ en cada paso:

\begin{equation}
    \hat{\mathbf{K}} = \tilde{\mathbf{K}}^{\top} \, \mathbf{P}_{N[:T,:]} \in \mathbb{R}^{B \times H \times d_c \times k}
    \label{eq:k_comprimido}
\end{equation}

El resultado es una representación comprimida de las claves en un espacio latente de
$k \ll T$ vectores base, cuyo coste de almacenamiento es $\mathcal{O}(T \cdot k)$ en
lugar de $\mathcal{O}(T^2)$.

\subsection{Cálculo de atención aproximada de bajo rango}
\label{subsec:cheap_attn}

A partir de las representaciones comprimidas, el predictor calcula una atención
aproximada (\textit{cheap attention}) mediante la siguiente operación:

\begin{equation}
    \mathbf{A}_{cheap} = \text{softmax}\!\left(\tilde{\mathbf{Q}} \cdot \hat{\mathbf{K}} \cdot s\right)
    \in \mathbb{R}^{B \times H \times T \times k}
    \label{eq:cheap_attn}
\end{equation}

donde $s = H^{-1/2}$ es el factor de escala estándar. Los coeficientes resultantes
$\mathbf{A}_{cheap}$ codifican, para cada posición temporal, el grado de afinidad con
cada uno de los $k$ vectores base del espacio comprimido. Estos coeficientes se utilizan
para reconstruir una aproximación de la matriz de atención completa $N \times N$ mediante
la proyección inversa $\mathbf{P}_{N}^{\top}$:

\begin{equation}
    \hat{\mathbf{A}} = \mathbf{A}_{cheap} \cdot \mathbf{P}_{back}^{\top}
    \in \mathbb{R}^{B \times H \times T \times T}
    \label{eq:approx_attn}
\end{equation}

donde $\mathbf{P}_{back} \in \mathbb{R}^{T_{max} \times k}$ es una segunda matriz de
proyección aprendida independiente de $\mathbf{P}_N$ (\texttt{share\_inout\_proj=False}),
que permite una mayor flexibilidad en la reconstrucción.

\subsection{Generación de la máscara binaria de atención}
\label{subsec:mascara_binaria}

La matriz de atención aproximada $\hat{\mathbf{A}}$ se utiliza para determinar qué
conexiones merecen ser calculadas con atención completa. Dado un presupuesto de atención
$b = \lceil r_{attn} \cdot T \rceil$, donde $r_{attn} = 0.25$ es la tasa de retención
definida en \texttt{transformer.py}, se seleccionan para cada posición de consulta los
$b$ índices de clave con mayor puntuación aproximada mediante una operación \textit{top-k}:

\begin{equation}
    \mathcal{I}^{(i)} = \underset{j \in [T]}{\text{top-}b}\;\hat{A}_{ij}
    \qquad \forall\, i \in [T]
    \label{eq:topk_selection}
\end{equation}

A partir de estos índices se construye la máscara binaria de atención
$\mathbf{M}_{attn} \in \{0,1\}^{B \times H \times T \times T}$, donde
$M_{attn}^{(i,j)} = 1$ si y solo si $j \in \mathcal{I}^{(i)}$. Esta máscara se
desconecta del grafo de gradientes (\texttt{.detach\_()}) para que la operación
\textit{top-k} —no diferenciable— no bloquee la retropropagación. El coste de calcular
la máscara es $\mathcal{O}(T \cdot k)$, y el coste de aplicarla en la atención completa
se reduce de $\mathcal{O}(T^2)$ a $\mathcal{O}(T \cdot b) = \mathcal{O}(r_{attn} \cdot T^2)$,
con $r_{attn} = 0.25$ en este sistema.

Adicionalmente, la implementación introduce un ajuste dinámico del presupuesto
(\texttt{real\_k = min(self.attn\_budget, N)}), que garantiza la compatibilidad con
secuencias más cortas que la longitud máxima $T_{max}$ sin requerir relleno hasta dicha
longitud. Esta es una contribución de ingeniería relevante respecto a la implementación
original de Sparsifiner, que asumía una longitud de secuencia fija.

\subsection{Atención completa con máscara dispersa}
\label{subsec:attn_completa_dispersa}

Una vez obtenida la máscara $\mathbf{M}_{attn}$, la capa de atención (\texttt{Attention})
calcula la matriz de atención estándar sobre las representaciones originales no comprimidas,
pero fuerza a $-\infty$ las posiciones no seleccionadas por la máscara antes de aplicar
\textit{softmax}:

\begin{equation}
    \mathbf{A} = \text{softmax}\!\left(\frac{\mathbf{Q}\mathbf{K}^{\top}}{\sqrt{d_h}}
    + \log \mathbf{M}_{attn}\right)
    \in \mathbb{R}^{B \times H \times T \times T}
    \label{eq:attn_dispersa}
\end{equation}

donde la suma con $\log \mathbf{M}_{attn}$ implementa el enmascaramiento: las posiciones
con $M_{attn}^{(i,j)} = 0$ reciben una puntuación de $-\infty$, resultando en un peso de
atención efectivo de cero tras el \textit{softmax}. Las filas de la matriz resultante que
corresponden a posiciones completamente enmascaradas —situación posible en secuencias
con \textit{padding}— son gestionadas mediante \texttt{torch.nan\_to\_num} para
garantizar la estabilidad numérica. La salida del módulo de atención es finalmente:

\begin{equation}
    \mathbf{Z}_{attn} = \mathbf{A} \mathbf{V} \mathbf{W}_O
    \in \mathbb{R}^{B \times T \times d}
    \label{eq:salida_attn}
\end{equation}

donde $\mathbf{V}$ son las representaciones de valor y $\mathbf{W}_O \in \mathbb{R}^{d \times d}$
es la proyección de salida.

% -----------------------------------------------------------
\section{Decodificador autorregresivo}
\label{sec:decoder}
% -----------------------------------------------------------

El decodificador del sistema, implementado en la clase \texttt{MidiDecoder}, adopta la
arquitectura estándar del decodificador Transformer \cite{vaswani2017attention}, compuesta
por $L_{dec} = 4$ capas apiladas de tipo \texttt{nn.TransformerDecoderLayer}. A diferencia
del codificador, el decodificador emplea atención densa estándar, decisión justificada por
las diferencias fundamentales entre las dos secuencias que procesa: mientras la secuencia
de audio puede contener decenas de miles de tramas, la secuencia de \textit{tokens} MIDI
generada tiene una longitud típicamente mucho menor, lo que hace manejable el coste
cuadrático de la autoatención.

\subsection{Embedding de entrada y codificación posicional}

Cada \textit{token} MIDI de la secuencia objetivo $\mathbf{y} \in \mathbb{Z}^{B \times S}$
se transforma en un vector de embedding de dimensión $d = 256$ mediante una tabla de
embeddings aprendida $\mathbf{E}_{tok} \in \mathbb{R}^{|\mathcal{V}| \times d}$, con el
índice de \textit{padding} (índice 0) inicializado a cero y excluido de la actualización
de gradientes. El embedding resultante se escala por $\sqrt{d}$ para compensar la
magnitud del embedding posicional \cite{vaswani2017attention}, y se le suma una
codificación posicional sinusoidal fija generada mediante la clase
\texttt{PositionalEncoding}:

\begin{align}
    PE_{(pos, 2i)}   &= \sin\!\left(\frac{pos}{10000^{2i/d}}\right) \label{eq:pe_sin} \\
    PE_{(pos, 2i+1)} &= \cos\!\left(\frac{pos}{10000^{2i/d}}\right) \label{eq:pe_cos}
\end{align}

La implementación utiliza la formulación logarítmicamente estable
$\exp\!\left(2i \cdot \frac{-\log(10000)}{d}\right)$ para el cálculo del denominador,
y el \textit{buffer} posicional se registra como parámetro no entrenable con longitud
máxima $max\_len = 50\,000$ \textit{tokens}, valor suficientemente holgado para cubrir
las secuencias MIDI más largas del corpus MAESTRO. La codificación posicional sinusoidal
se ha elegido para el decodificador —en contraposición al embedding posicional aprendido
del codificador— por su capacidad de generalización a longitudes de secuencia no vistas
durante el entrenamiento \cite{vaswani2017attention}.

\subsection{Estructura de cada capa del decodificador}

Cada capa del decodificador implementa el flujo de procesamiento canónico en tres
subcapas:

\paragraph{Autoatención enmascarada (\textit{Masked Self-Attention}).}
La primera subcapa aplica autoatención multi-cabeza sobre la secuencia de \textit{tokens}
ya generados, utilizando una máscara triangular inferior (\textit{causal mask}) para
garantizar la propiedad autorregresiva: la predicción del \textit{token} en la posición
$t$ solo puede depender de los \textit{tokens} en posiciones $0, 1, \ldots, t-1$. La
máscara se genera dinámicamente en cada paso de decodificación mediante:

\begin{equation}
    M_{causal}^{(i,j)} = \begin{cases} 0 & \text{si } j \leq i \\ -\infty & \text{si } j > i \end{cases}
    \label{eq:mascara_causal}
\end{equation}

\paragraph{Atención cruzada (\textit{Cross-Attention}).}
La segunda subcapa aplica atención cruzada entre la secuencia de \textit{tokens} del
decodificador (actuando como consultas) y la memoria $\mathbf{M}$ producida por el
codificador (actuando como claves y valores). Este mecanismo es el que permite al
decodificador acceder al contenido acústico de la pieza en cada paso de generación,
condicionando la predicción del siguiente \textit{token} MIDI a la representación latente
completa del espectrograma de entrada.

\paragraph{Red \textit{feed-forward}.}
La tercera subcapa es una red de propagación hacia adelante de dos capas lineales con
función de activación ReLU y dimensión oculta $d_{ff} = 2048$, seguida de \textit{dropout}
con $p = 0.1$.

Todas las subcapas implementan conexiones residuales y normalización de capa
(\textit{Post-LayerNorm}), siguiendo la implementación estándar de PyTorch \newline
(\texttt{nn.TransformerDecoderLayer} con \textit{batch\_first=True}).

\subsection{Capa de clasificación y generación}

La salida de la última capa del decodificador,
$\mathbf{H} \in \mathbb{R}^{B \times S \times d}$, se proyecta a logits sobre el
vocabulario mediante una capa lineal de salida:

\begin{equation}
    \mathbf{L} = \mathbf{H} \, \mathbf{W}_{out} + \mathbf{b}_{out}
    \in \mathbb{R}^{B \times S \times |\mathcal{V}|}
    \label{eq:logits}
\end{equation}

donde $\mathbf{W}_{out} \in \mathbb{R}^{d \times |\mathcal{V}|}$ con $|\mathcal{V}| = 311$ (véase Sección~\ref{sec:dominio_simbolico}).
Durante el entrenamiento, estos logits se utilizan directamente para calcular la función
de pérdida de entropía cruzada con la secuencia objetivo. Durante el entrenamiento se emplea \textit{teacher forcing}
estricto: el decodificador recibe como entrada la secuencia de
\textit{tokens} objetivo desplazada una posición
(\texttt{batch\_midi[:, :-1]}), lo que acelera la convergencia
pero introduce una discrepancia entre las condiciones de
entrenamiento e inferencia (\textit{exposure bias}).

Durante la inferencia el sistema soporta dos estrategias de
decodificación. La primera es la decodificación voraz
(\textit{greedy decoding}): en cada paso $t$ se selecciona el
\textit{token} con mayor logit y se incorpora a la secuencia de
entrada para el paso siguiente. La segunda es el muestreo con
temperatura (\textit{temperature sampling}), con
$\mathit{temperature} = 0.8$, que introduce estocasticidad en
la selección mediante \texttt{torch.multinomial}, favoreciendo
la diversidad de las transcripciones generadas. En ambos casos
el proceso se detiene al predecir el \textit{token}
\texttt{<eos>} (índice~310) o al alcanzar la longitud máxima
de generación.

% -----------------------------------------------------------
\section{Resumen de la arquitectura y flujo de datos}
\label{sec:resumen_arquitectura}
% -----------------------------------------------------------

El flujo de datos completo a través del modelo fue introducido esquemáticamente en la Figura~\ref{fig:encoder_decoder} del
Capítulo~\ref{cap:estado_del_arte}. A nivel de arquitectura global,
el codificador \newline\texttt{AudioSparsifinerEncoder} ocupa el bloque
izquierdo del diagrama, produciendo la memoria $\mathbf{M}$ que
el decodificador consume mediante atención cruzada. El mecanismo
interno del codificador —proyección de entrada, bloques Sparsifiner
apilados y predictor de máscara— queda descrito con precisión en
las ecuaciones de las Secciones~\ref{sec:encoder_sparsifiner}
y~\ref{sec:mask_predictor}. En síntesis, la arquitectura \texttt{Magdalena v1.0}
se caracteriza por los siguientes aspectos diferenciales respecto a un Transformer
estándar:

\begin{itemize}
    \item \textbf{Codificador asimétrico con atención dispersa:} El codificador aplica
    atención dispersa instancia-dependiente mediante el mecanismo Sparsifiner, reduciendo
    el coste de atención efectivo de $\mathcal{O}(T^2)$ a $\mathcal{O}(r_{attn} \cdot T^2)$
    con $r_{attn} = 0.25$, sin comprometer la capacidad del modelo para capturar
    dependencias globales seleccionadas dinámicamente.

    \item \textbf{Predicción de máscara de bajo rango:} La selección de conexiones
    relevantes se realiza mediante una factorización de bajo rango de la matriz de atención
    con factores de compresión $r_c = 2$ (canal) y $r_N = 64$ (temporal), lo que mantiene
    el coste del predictor en $\mathcal{O}(T \cdot k)$ con $k = T / r_N \ll T$.

    \item \textbf{Compatibilidad con longitud variable:} La implementación introduce
    recorte dinámico (\textit{slicing}) de las matrices de proyección del predictor,
    eliminando la restricción de longitud fija del Sparsifiner original y permitiendo
    procesar fragmentos de audio de duración arbitraria sin necesidad de \textit{padding}
    hasta $T_{max}$.

    \item \textbf{Decodificador estándar con atención cruzada:} El decodificador emplea
    la arquitectura Transformer canónica con autoatención enmascarada y atención cruzada
    hacia la memoria del codificador, garantizando la coherencia gramatical de la
    secuencia de \textit{tokens} MIDI generada mediante la máscara causal.
\end{itemize}

En conjunto, estas decisiones de diseño permiten que
\texttt{Magdalena v1.0} procese secuencias de audio de hasta
82 segundos con los recursos de cómputo disponibles, algo
inviable con un Transformer estándar a igual configuración.
La reducción de dimensionalidad respecto a la arquitectura
de referencia ($d = 256$ frente a $d = 512$) es una concesión
a las limitaciones de hardware que el mecanismo Sparsifiner
compensa parcialmente al reducir el coste cuadrático de la
atención. El Capítulo~\ref{cap:implementacion} describe los
detalles de entrenamiento y la estrategia experimental adoptada
para validar estas elecciones dentro de las restricciones
existentes.
